\chapter{Théorie de la d\'ecision}

\section{Introduction}
La théorie de la décision modélise le comportement d'un agent face
à des situations de choix. Nous décrivons les choix d'un agent, et
les relions à une notion de préférences. Nous montrons sous
quelles conditions les relations de pr\'ef\'erences peuvent \^etre
repr\'esent\'ees par une fonction d'utilit\'e, et quand elles
impliquent l'existence de croyances de l'agent sur les
diff\'erents \'etats du monde, repr\'esentables par des
probabilit\'es.

\begin{comment}Je crois qu'il serait bien de parler des deux aspects :
normatif et descriptif, car cela aide bien à comprendre le cours
par la suite, et est vrai pour toute la théorie des jeux et la
théorie économique en general.
\end{comment}

\section{Choix et préférences}
Nous commençons par décrire des préférences d'un agent, et des
choix, et établissons un lien naturel entre ces deux notions.
\subsection{Choix}
Nous partons d'un ensemble fini $X$ d'alternatives. Un agent peut
effectuer un choix parmi un sous-ensemble de $X$ d'alternatives
qui lui sont proposées.  Soit donc $2^X$ l'ensemble des parties
de~$X$.
\begin{definition}
Une fonction de choix est la donnée de $c\colon
2^X\backslash\{\emptyset\}\to 2^X\backslash\{\emptyset\}$ telle
que pour tout $A\subset X$, $c(A)\subset A$.
\end{definition}
L'interprétation est la suivante : si l'agent a le choix entre les
éléments de $A$, cet agent dira que tous les éléments de $c(A)$
conviennent, ou bien choisira indifféremment parmi les éléments de
$c(A)$.

 Certaines fonctions de choix vérifient des propriétés qui les
rendent plus ``plausibles'' que d'autres. Nous étudierons plus
particulièrement l'influence des axiomes suivant, appelés axiomes
$(\alpha)$ et $(\beta)$ de Sen :
\begin{itemize}
\item{\bf $(\alpha)$:} Si $x\in B\subset A$ et $x\in c(A)$, alors
$x\in c(B)$ %
\item{\bf $(\beta)$:} Si $x,y\in c(A)$, $A\subseteq B$, et $y\in
c(B)$, alors $x\in c(B)$.
\end{itemize}
On peut traduire ces axiomes en langage courant de la manière
suivante: Pour l'axiome $(\alpha)$, ``Si le champion du monde d'un
jeu est pakistanais, alors ce pakistanais doit aussi \^etre
champion du Pakistan.'' Pour l'axiome $(\beta)$, ``Si le champion
du monde est pakistanais, alors tous les champions du Pakistan
sont aussi champions du monde.''

\begin{comment}
En cours, donner des exemples de fonctions de choix, et montrer
qu'elles vérifient (ou non) certains des axiomes.
\end{comment}

\subsection{Préférences}

\begin{definition}
Une relation de préférences $\succ$ est la donnée d'un ensemble de
paires d'éléments $(x,y)$ de $X\times X$ pour lesquels nous
écrivons $x\succ y$. La relation $x\succ y$ est interprétée par :
``l'agent préf\`ere strictement $x$ à~$y$''.
\end{definition}

Voici une liste d'axiomes sur les relations de préférences:
\begin{itemize}
\item[\bf Asymétrie] $x\succ y\implies \neg y\succ x$. Si $x$ est
strictement préféré à $y$, alors $y$ n'est pas strictement préféré
à $x$. %
\item[\bf Transitivité] $(x\succ y\mbox{ et }y\succ z)\implies
x\succ z$. %
\item[\bf Irreflexivité] $\neg x\succ x$. Il n'existe pas de $x$
strictement préféré à lui-m\^eme.%
\item[\bf Transitivité négative] $(\neg x\succ y \mbox{ et }\neg
y\succ z)\implies \neg x\succ z$.%
\item[\bf Acyclicité] $x_1\succ x_2\succ\ldots\succ x_n\implies
\neg x_n\succ x_1$. Si chaque $x_i$ est strictement préféré à
$x_{i+1}$, alors $x_n$ n'est pas préféré à $x_1$.
\end{itemize}

\begin{definition}
Nous dirons qu'une relation de préférences est rationnelle si elle
est asymétrique et négativement transitive.
\end{definition}
Une relation de préférences ainsi définie possède forcément les
autres propriétés de la liste:
\begin{proposition}\label{rationelleimpliquelereste}
Une relation de préférences rationnelle est irreflexive,
transitive, et acyclique.
\end{proposition}
Une relation de préférences rationnelle est un ordre partiel.

A partir d'une relation binaire $\succ$, nous pouvons définir une
relation de préférences faibles $\succeq$  et une relation
d'indifférence $\sim$ par:
\begin{eqnarray*}
x\succeq y &\mbox{\ si\ }& \neg y\succ x \\
x\sim y &\mbox{\ si\ }& \neg y\succ x\mbox{\ et\ } \neg x\succ y %
\end{eqnarray*}

Voici quelques propriété des relations $\succeq$ et $\sim$ lorsque
$\succ$ est rationnelle.

\begin{proposition}
Supposons que $\succ$ est une relation de préférences rationnelle,
alors
\begin{itemize}
\item Pour tout $x$ et $y$, on a soit $x\succ y$, soit $y\succ x$,
soit $x\sim y$.%
\item $x\geq y$ est compl\^ete (on a toujours $x\geq y$ ou $y\geq
x$) et transitive. %
\item $\sim$ reflexive ($x\sim x$ pour tout $x$), symétrique
($x\sim y$ si et seulement si $y\sim x$), et transitive. %
\item $w\succ x$, $x\sim y$ et $y\succ z$ impliquent $w\succ y$ et
$x\succ z$.%
\item $x\geq y$ si et seulement si $x\succ y$ ou $x\sim y$. %
\item $x\geq y$ et $y\geq x$ si et seulement si $x\sim y$.
\end{itemize}
\end{proposition}

\subsection{Lien entre choix et préférences}
Deux idées simples et importantes relient les choix et les
préférences. Si un agent possède des préférences, ceci doit avoir
une implication en termes de choix. Lorsqu'on cherche quelles sont
les relations de préférences compatibles avec des choix (s'il en
existe), on parle de préférences révélées.

\subsubsection{Comment les préférences définissent des choix}
Soit $\succ$ une relation binaire sur $X$. On définit une fonction
$c(\cdot,\succ)$ par :
$$
c(A,\succ) = \{x\in A \colon \mbox{ pour tout $y\in A$ } \neq y \succ x\}
$$
$c(A,\succ)$ est donc le sous-ensemble des éléments de $A$ qui ne
sont pas dominés par d'autres éléments de $A$. Si $\succ$
représente les préférences de l'agents, alors ses choix doivent
obéir a $c(\cdot,\succ)$. Sous quelles conditions sommes-nous
s\^urs que $c(\cdot,\succ)$ est une fonction de choix ?

\begin{proposition}
Étant donnée une relation binaire $\succ$, $c(\cdot,\succ)$ est
une fonction de choix si et seulement si $\succ$ est acyclique.
\end{proposition}

Montrons que l'acyclicité implique $c(B,\succ)\neq\emptyset$.
C'est une consequence de $X$ fini. Réciproquement, $\succ$ n'est
pas acyclique, soit $x_1\succ x_2\succ\ldots\succ x_n\succ x_1$.
On a alors $c(\cdot, \succ)(\{x_1,x_2,\ldots, x_n\})=\emptyset$.

\subsubsection{Représentation des choix par des préférences rationnelles}
Nous partons maintenant d'une fonction de choix $c$, et cherchons
sous quelles conditions il existe une relation de préférences
rationnelle $\succ$ telle que $c$ soit la fonction de choix
correspondant à $\succ$ ($c(\cdot,\succ)=c$).

\begin{proposition}
Étant donnée une fonction de choix $c$, il existe une relation
rationnelle $\succ$ telle que $c=c(\cdot,\succ)$ si et seulement
si $\succ$ vérifie les axiomes $(\alpha)$ et $(\beta)$ de Sen. De
plus, cette relation est alors unique.
\end{proposition}

 Démonstration : Tout d'abord, supposons $c=c(\cdot,\succ)$, avec $\succ$
rationnelle. Pour montrer $\alpha$, il suffit de remarquer que
$\not\exists y \in A, y\succ x$ implique $\not\exists y \in B,
y\succ x$, pour $B$ sous-ensemble de $A$. Pour $\beta$, supposons
$(x,y)\in c(A)$, on a alors $x\sim y$. Si $x\in c(B)$ avec
$A\subset B$, $\forall z\in B$, $z\preceq x$, et donc $z\preceq y$
par transitivité, donc $y\in c(B)$.\\
 Supposons maintenant que $c$ vérifie $\alpha$ et $\beta$.
L'unicité de $\succ$ provient de ce que $\succ$ doit vérifier
$x\succ y$ si et seulement si $c(\{x,y\})= \{x\}$. La reflexivité
de $\succ$ est triviale. Pour la transitivité de $\succeq$,
supposons $x\succeq y\succeq z$. Si $y\in c(\{x,y,z\})$ alors
$z\in c(\{x,y,z\}$ ($\beta$). Sinon, on ne peut avoir
$c(\{x,y,z\}=\{x\}$ (car alors on aurait $c(\{x,y,z\}=\{x\}$).
Donc $z\in c(\{x,y,z\}=\{x\}$. On a donc aussi $z\in c(\{x,z\}$
($\alpha$), et $x\succeq z$. La relation $\succ$ est donc
rationnelle. Il reste à montrer que $c=c(\cdot, \succ)$, c'est à
dire que $x\in c(A)$ si et seulement si pour tout $y\in A$, $x\in
c(\{x,y\})$. Pour la partie ``seulement si'', $\forall y$, $x\in
c(\{x,y\})$, par l'axiome $\alpha$. Pour la partie ``si'', soit
$z\in c(A)$, alors $z\in c(\{x,z\})$ ($\alpha$), donc $c(\{x,z\})=
\{x,z\}$, et ($\beta$) implique $x\in c(A)$.


Il suffit d'observer les choix faits par les agents ($c$) pour en
déduire leurs préférences ($\succ$). Pour cette raison, la
conjonction des axiomes ($\alpha$) et ($\beta$) est aussi appelée
{\em axiome des préférences révélées}.

\section{Représentation ordinale}
Soit $u$ une fonction sur $X$. On interprète $u(x)$ comme une
utilité, ou plaisir, que l'agent peut retirer de l'alternative
$x$. On appelle donc $u$ {\em fonction d'utilité}. Si $u$ est la
fonction d'utilité de l'agent, quelles doivent \^etre ses
préférences ? Simplement, définissons $\succ_u$ par
$$
x\succ_u y \mbox{ si et seulement si } u(x) > u(y)
$$
Il est aisé de vérifier que $\succ_u$ est bien une relation de
préférences rationnelle. La proposition suivante nous dit aussi
que toute relation de préférences résulte d'une fonction d'utilité
($X$ est fini).
\begin{proposition}\label{repsiXfini}
Supposons que $X$ est fini, et soit $\succ$ une relation binaire
sur $X$. Il existe une fonction d'utilité $u$ telle que $\succ_u =
\succ$ si et seulement si $\succ$ est rationnelle.
\end{proposition}
Preuve : par induction sur le cardinal de $X$.

La proposition \ref{repsiXfini} s'étend au cas $X$ dénombrable.

\begin{comment}
Pourquoi ne pouvons pas faire d'induction transfinie ? Quelles
seraient les axiomes permettant une induction transfinie ?
Conjecture: probablement qqch + fort que l'acyclicité, genre pas
de cycles infinis...
\end{comment}

\begin{proposition}\label{repsiXdenom}
Supposons $X$ dénombrable, et soit $\succ$ une relation binaire
sur $X$. Il existe une fonction d'utilité $u$ telle que $\succ_u =
\succ$ si et seulement si $\succ$ est rationnelle.
\end{proposition}

Nous pouvons nous poser la question de l'unicité de la
représentation d'une relation de préférences par une fonction
d'utilité. Il est évident que la composition de $u$ par une
fonction strictement croissante ne change par $\succ_u$. La
proposition suivante nous donne une réciproque:

\begin{proposition}\label{unicitepreprordinale}
Soient $u$ et $u'$ deux fonctions d'utilité sur $X$ (fini). Alors
$\succ_u=\succ_{u'}$ si et seulement si il existe une fonction
$f\colon \reals \to \reals$ telle que:
\begin{itemize}
\item $f$ est strictement croissante%
\item $u'=f\circ u$
\end{itemize}
\end{proposition}
On commence par définir $f$ sur $\{u(x), x\in X\}$. Puis on
l'étend à $\reals$.

\section{Préférences sur choix incertains}
Introduction. Expliquer la différence entre choix incertains dont
on peut évaluer les probabilités de manière objective ou non.

\subsection{Préférences sur les probabilités objectives}

Soit $Z$ un ensemble fini, appelé ensemble de prix. Soit
$P=\Delta(Z)$ l'ensemble des lois de probabilités sur $Z$, appelé
ensemble des loteries. Nous étudions maintenant des relations de
préférences sur $P$. $P$  est donc l'ensemble de choix.

 Il est remarquable que $P$ a une structure particulière, appelée
structure de simplexe. Nommément :
\begin{itemize}
\item $P$ est convexe. Pour deux éléments $p$ et $q$ de $P$, et
$\lambda\in[0,1]$, $\lambda p + (1-\lambda) q \in P$, o\`u
$(\lambda p + (1-\lambda) q)(x)$ vaut par définition $\lambda p(x)
+ (1-\lambda) q(x)$. %
\item On peut identifier chaque $z\in Z$ à l'élément de $P$ qui
donne une probabilité de 1 à $z$, et $0$ aux autres.
\end{itemize}

 On définit maintenant des axiomes sur la relation de préférences
$\succ$ de $P$ qui utilisent la structure particulière de $P$.

\begin{itemize}
\item[{\bf Rationalité}] $\succ$ est une relation de
préférences rationnelle.%
\item[\bf Indépendance] Pour $p,q,r\in P$, et $0<\lambda\leq 1$,
$p\succ q$ implique $\lambda p+(1-\lambda)r\succ \lambda
q+(1-\lambda)r$%
\item[\bf Continuité] Pour $q\in P$, les ensembles $\{p, p\succ
q\}$ et $\{p, q\succ p\}$ sont ouverts
\end{itemize}

L'axiome de continuité se comprend lorsqu'on se rappelle que
$\succ$ représente une relation de préférences stricte. Si $p\succ
q$, alors si on modifie de manière extr\^emement mince $p$ en $p'$
et $q$ en $q'$ on aura $p'\succ q'$.

\subsection{Utilité espérée}
Supposons que l'agent associe une valeur, ou utilité $u(z)$ à
chacun des prix $z\in Z$. La fonction $u$ définit une fonction
$\tilde u$ de $P$ vers $\reals$ donnée par:
$$
\tilde u(p) = \sum_z p(z) u(z)
$$
Soit $\succ_{\tilde u}$ la relation de préférences sur $P$
associée à $\tilde u$.

Remarquons que nous n'avons aucune raison de supposer {\em a
priori} que les préférences de l'agent soient de la forme
$\succ_{\tilde u}$, pour un $u$ donné, ni encore qu'il existe $u$
telle que $\succ = \succ_{\tilde u}$. Il peut d'ailleurs sembler à
première vue que cette forme de préférences est restrictive et
très particulière. Cependant, cette forme de préférences est très
commode pour le modélisateur (nous), ou pour un agent cherchant à
évaluer ses préférences entre des loteries, car la donnée
d'utilités associées aux événements certains (les prix de $Z$) à
elle seule définit les préférences sur l'ensemble beaucoup plus
complexe des loteries.

\subsection{Théorème de von-Neumann et Morgenstern}
Le résultat de von-Neumann et Morgenstern est l'équivalence entre
une famille d'axiomes sur $\succ$ et l'existence d'une
représentation de la forme $\succ_{\tilde u}$.
%
\begin{proposition}
Il existe une fonction $u:Z\to \reals$ telle que $\succ =
\succ_{\tilde u}$ si et seulement si $\succ$ vérifie les axiomes
de préférences, d'indépendance et de continuité.
\end{proposition}
%
{\bf Démonstration de la partie ``si''} Il suffit de montrer que
$\succ_{\tilde u}$ vérifie les propriétés requises. \\
{\bf Démonstration de la partie ``seulement si''} Pour $z\in Z$,
on note $\delta_z$ l'élément de $P$ qui met probabilité 1 sur $Z$.
Si pour tous $x,y\in Z$, $\delta_x\sim\delta_y$, alors pour tous
$p,p'\in P$, $p\sim p'$ (indépendance), et toute fonction $u$
constante convient. Dans le cas contraire, on sélectionne parmi
l'ensemble $\{\delta_z,\ z\in Z\}$ un élément minimal
$\delta_{x_0}$, et un élément maximal $\delta_{x_1}$ pour $\succ$.
D'après l'axiome de continuité, pour chaque $z$, il existe un
unique $\lambda_z\in[0,1]$ tel que $z\sim\lambda_z
\delta_{x_1}+(1-\lambda_z)\delta_{x_0}$. On pose alors
$u(z)=\lambda_z$. On a donc $u(x_0)=0$, $u(x_1)=1$, et pour $p\in
P$ $\tilde u(p)= \sum_k p(z)\lambda(z)$. Soient $p$ et
$p'$ dans $P$.%
\begin{eqnarray*}
p\succ p' & \Leftrightarrow & \sum_z p'(z)(\lambda_z
\delta_{x_1}+(1-\lambda_z)\delta_{x_0}) \succ \sum_z
p'(z)(\lambda_z \delta_{x_1}+(1-\lambda_z)\delta_{x_0})\\
& \Leftrightarrow & (\sum_z p(z)\lambda(z))\delta_{x_1}+(\sum_z
p(z)(1-\lambda(z)))\delta_{x_0} \succ \\
&&(\sum_z p(z)'\lambda(z))\delta_{x_1}+(\sum_z
p(z)'(1-\lambda(z)))\delta_{x_0}\\
& \Leftrightarrow & \sum_z  p(z)\lambda(z) > \sum_z p'(z)\lambda(z)\\
& \Leftrightarrow & \tilde u(p) > \tilde u(p')
\end{eqnarray*}
\begin{proposition}
Deux fonctions d'utilité $\tilde u$ et $\tilde u'$ sur $P$ sont
associées aux m\^emes préférences si et seulement il existe $a>0$
et $b$ tels que $\tilde u' = a\tilde u+b$.
\end{proposition}
\noindent
{\bf Démonstration de la partie ``si''} est évidente.\\
{\bf Démonstration de la partie ``seulement si''} Supposons que
$\succ_{\tilde u} = \succ_{\tilde u'}$, et soient $\delta_{x_0}$,
$\delta_{x_1}$ un élément minimal et un élément maximal. Pour tout
$p\in P$, écrivons $\tilde u(p)= \frac{u(x_1)-\tilde
u(p)}{u(x_1)-u(x_0)}u(x_0)+\frac{\tilde
u(p)-u(x_0)}{u(x_1)-u(x_0)}u(x_1)$. Donc $p\sim_{\tilde
u}\frac{u(x_1)-\tilde
u(p)}{u(x_1)-u(x_0)}\delta_{x_0}+\frac{\tilde
u(p)-u(x_0)}{u(x_1)-u(x_0)}\delta_{x_1}$. Comme $\sim_{\tilde
u}=\sim_{\tilde u'}$,
$$%
\frac{u(x_1)-\tilde u(p)}{u(x_1)-u(x_0)} = \frac{u'(x_1)-\tilde
u(p)}{u'(x_1)-u'(x_0)}
$$%
Cad.
$$%
\tilde u'(p)= \frac{u'(x_1)-u'(x_0)}{u(x_1)-u(x_0)}u(p)-
\frac{u'(x_1)-u'(x_0)}{u(x_1)-u(x_0)}u(x_1) + u'(x_1)
$$%
ce qui est bien la forme recherchée.
\subsection{Préférences et incertain subjectif}
Dans une loterie, l'agent à un moyen objectif d'estimer la
probabilité de chaque prix. En revanche, dans un événement comme
une course de chevaux par exemple, il n'existe pas de moyen
objectif de conna\^itre la probabilité que tel ou tel cheval soit
vainqueur. La théorie de Anscombe et Aumann nous donne des
conditions sous lesquelles l'agent se comporte {\em comme si} il
possédait une croyance sur les chevaux vainqueurs, et effectuait
ses choix (ici des paris) en maximisant son utilité espérée étant
donnée cette croyance. Ces probabilités sont dites probabilités
subjectives.

\subsubsection{Paris sur des loteries}
$S$ est un ensemble fini de possibilités incertaines (comme les
chevaux gagnants dans le cas de la course). $Z$ est un ensemble
fini de prix, et $P$ représente les probabilités sur $Z$. A chaque
élément de $S$, faisons correspondre un prix aléatoire. Imaginons
que pour chaque élément de $S$, une loterie puisse \^etre tirée
pour déterminer le prix dans $Z$. On définit alors une {loterie
composée} comme un vecteur $h =(h_s)_{s\in S}$ de $P^S$.
L'ensemble de choix que nous considérons maintenant l'ensemble $H=
P^S$ des loteries composées.

Exemples de loteries composées :
\begin{itemize}
\item S'il pleut, mon cheval préféré Rossinante a 45\% de chances
de remporter la course, et seulement 20\% de chances s'il ne pleut
pas.%
\item Si l'économie est en croissance l'an prochain, j'ai une
chance sur deux d'obtenir une promotion, et une chance sur trois
sinon.
\end{itemize}

 Remarquons que $H$ a lui aussi une structure de simplexe. En
particulier, si $h,h'\in H$ et si $\lambda\in [0,1]$, $h''$ défini
par $h_s''(z)= \lambda h_s(z)+(1-\lambda)h_s'(z)$ appartient à
$H$.

\subsubsection{Les axiomes}
Supposons que l'agent a des préférences sur $H$, représentées par
$\succ$. Comme dans le cas des probabilités objectives, soient les
axiomes suivants sur $\succ$.

\begin{itemize}
\item[Axiome de rationalité] $\succ$ est une relation de
préférences rationnelle (sur $H$). \item[Axiome d'indépendance]
Pour $h,h',h''\in H$, et $0<\lambda\leq 1$, $h\succ h'$ implique
$\lambda h+(1-\lambda)h''\succ \lambda h+(1-\lambda)h'$
\item[Axiome de continuité] Pour $h,h',h''\in H$, si $h\succ
h'\succ h''$ il existe $0<\lambda,\mu<1$ tels que $\lambda
h+(1-\lambda)h''\succ h'\succ \mu h+(1-\mu)h''$
\end{itemize}

\subsubsection{Utilité dépendante de l'état}
Comme pour le cas des fonctions d'utilité de von-Neumann et
Morgenstern, nous pouvons caractériser les relations de binaires
qui vérifient les axiomes de  préférences, d'indépendance et de
continuité.
%
\begin{proposition}
Une relation binaire $\succ$ vérifie les axiomes de préférences,
d'indépendance et de continuité si et seulement si il existe une
famille $u=(u_s)_{s\in S}$ d'applications $u_s\colon Z\to \reals$
telles que
$$
h\succ h' \Leftrightarrow \sum_{s\in S}\sum_{z\in Z}
h_s(z)u_s(z)>\sum_{s\in S}\sum_{z\in Z} h'_s(z)u_s(z)
$$
\end{proposition}
\begin{comment}
Parler d'unicité ? En parler pour les fonctions de vN-M ?
\end{comment}
La démonstration est la m\^eme que celle du théorème de
von-Neumann et Morgenstern.
%
On peut alors poser $\tilde u_s(h_s)=\sum_z h_s(z)u_s(z)$, utilité
de von-Neumann et Morgenstern dans l'état du monde $s$.

\subsubsection{Utilité indépendante de l'état et probabilités subjectives}

 Nous allons introduire un axiome disant que les préférences ne
dépendent pas de l'état du monde. Soit donc $\succ_{\tilde u_s}$
la relation de préférences dans l'état $s$.
\begin{itemize}
\item[\bf Indépendance de l'état] $\succ_{\tilde u_s}$ ne dépend
pas de $s$.
\end{itemize}
Selon l'axiome d'indépendance de l'état, si $p$ est préféré à $p'$
dans un état $s$, il est aussi préféré à $p'$ dans tous les états
$s'$.

Dans ce cas, fixons un état $s_0$, et posons $\tilde u = \tilde
u_{s_0}$. On a pour tout état $s$ l'existence de réels
$a_{s}>0,b_s$ tels que , $\tilde u_{s}=
a_{s}\tilde u_{s_0} +b_{s}$. On peut donc écrire:%
$$
\sum_{s} \tilde u_s(h_s) = \sum_s a_s \tilde u_{s_0}(h_s) +\sum_s b_s
$$
 On a donc $h\succ h'$ si et seulement si:
$$
\sum_s a_s \tilde u_{s_0}(h_s) > \sum_s a_s \tilde u_{s_0}(h'_s)
$$
Posons alors $\pi(s) = \frac{a_s}{\sum_s a_s}$. Tous les $\pi(s)$
sont positifs, et leur somme vaut 1. On peut donc interpréter
$\pi=(\pi(s))_s$ comme une probabilité sur $S$. On a alors $h\succ
h'$ si et seulement si:
$$
\sum_s \pi(s) \tilde u_{s_0}(h_s) > \sum_s \pi(s) \tilde
u_{s_0}(h'_s)
$$
Tout se passe donc comme si les préférences de l'agent étaient
représentées par les préférences $\tilde u$, et par une
probabilité $\pi$ sur $S$. Il est naturel d'interpréter $\pi$
comme les {\em croyances} de l'agent sur $S$. Ces croyances ne
sont pas une donnée objective, on parle de croyances {\em
subjectives} de l'agent.
%
\section{Exercices}
\subsection{Exercice 1}
Soit l'axiome suivant, appelé Axiome de Houthakker, portant sur
une fonction de choix $c$:\\
Si $x,y\in A\cap B$, et si $x\in c(A)$ et $y\in c(B)$, alors $x\in
c(B)$.\\
Montrer que l'axiome de Houthakker est équivalent à la conjonction
des axiomes $\alpha$ et $\beta$ de Sen.
\subsection{Exercice 2}
Montrer la proposition \ref{rationelleimpliquelereste} %
\subsection{Exercice 3}
Une relation de préférences non représentable par une fonction
d'utilité. Sur $X=[0,1]\times[0,1]$, on définit la relation de
préférences entre $x=(x_1,x_2)$ et $y=(y_1,y_2)$ par $x\succ y$ si
et seulement si $x_1> y_1$, ou ($x_1=y_1$ et $x_2>y_2$). C'est
donc l'ordre lexicographique sur $X$.
\begin{enumerate}
\item Montrer que $\succ$ est rationnelle.%
\item Montrer que $\succ$ n'est représentable par aucune fonction
d'utilité. C'est à dire, montrer qu'il n'existe pas d'application
$u$ de $X$ vers $\reals$ telle que $x\succ y$ si et seulement si
$u(x)> u(y)$.
\end{enumerate}
\subsection{Exercice 5} On a deux actifs financiers. Le premier
paye 100 dollars dans l'état du monde 1 (reprise de l'économie
américaine), et 50 dollars dans l'état du monde 2 (pas de reprise
de l'économie américaine). Le deuxième paye 120 dollars dans
l'état du monde 1 et 20 dans l'état du monde 2. Sur les marchés,
le prix du premier actif financier est de $80$ euros, celui du
second de 70 euros. Dans le ``Financial Times'' vous lisez le
commentaire suivant, ``d'après le marché, l'état du monde 1 et
l'état du monde 2 ont les mêmes chances de se réaliser.''
\begin{enumerate}
\item Ce commentaire est-il en rapport avec les prix que vous
observez ? %
\item  Selon vous, quelles sont les probabilités subjectives du
marché sur l'état 1 et sur l'état 2 ? %
\end{enumerate}
\subsection{Exercice 6: Paradoxe d'Allais}
On a comme ensemble d'alternatives $X =\{ 2 500 000 $\euro$, 500
000 $\euro$, 0 $\euro$ \}$, où $2 500 000 $\euro correspond au
fait de gagner cette somme, etc.\ldots On a les loteries $p_1 =
(0,1,0)$,  $p'_1 = (0.1,0.89,0.01)$, $p_2 = (0,0.11,0.89)$ et $p'2
= (0.1,0,0.9)$. (Par exemple, avec $p_2$, la probabilité de gagner
2500000\euro~ est de 0.11.
\begin{enumerate}
\item On suppose qu'un individu préfère $p_1$ à $p'_1$, et $p'_2$ à $p_2$.
Est-ce consistant avec l'utilité espérée de von-Neumann et Morgenstern ?
\end{enumerate}
\subsection{Exercice 7: Paradoxe de Machina}
$X = \{ $voyage à Venise, beau film sur Venise, rien $\}$, et $p_1
= (0.999,0.001,0)$, $p'_1 = (0.999,0,0.001)$. Si un agent préfère
voir un beau film sur Venise plut\^ot que rien, et préfère $p'_1$
à $p_1$, est-ce consistant avec la théorie de l'utilité espérée de
von-Neumann et Morgenstern ? Comment peut-on expliquer de telles
préférences ?
